
\subsubsection{2.1 Permutation Equivariant Weight Space Learning}

The observation that permuting neurons in adjacent layers preserves network function underlies a class of equivariant architectures for weight space learning. NFN [Zhou et al., NeurIPS 2023] formalizes this symmetry into equivariant layers for MLP and CNN weights, achieving Kendall's tau = 0.934 on CIFAR-10-GS and tau = 0.931 on SVHN-GS in within-architecture generalization prediction. Monomial-NFN [Tran, Vo et al., NeurIPS 2024] extends NFN to scaling and sign-flip symmetries, establishing that all invariant groups for MLPs are subsets of the monomial matrix group. Both methods remain architecture-specific: they are not designed to transfer representations across CNN and transformer families.

Transformer-NFN [Tran-Viet et al., 2024] identifies the maximal symmetric group of multi-head attention weights and reports within-architecture tau = 0.905-0.910 on transformer benchmarks. Critically, its equivariant construction for attention layers is separate from the CNN construction. The two constructions are not unified. The present work examines whether a single group — the input/output channel permutation group — can serve as a common foundation and measures what fraction of variance it captures for each layer type.

\subsubsection{2.2 Token-Based and Scalable Weight Representations}

SANE [Schurholt et al., ICML 2024] tokenizes neural network weights by reshaping tensors into row-wise slices along the output channel dimension and encoding each token with a sequential positional encoding P_n = [n, l, k]. SANE achieves linear probe accuracy of 0.978 on MNIST and 0.991 on CIFAR-10 within architecture, and scales to ResNets. Its sequential positional encodings carry no symmetry structure, which is hypothesized to limit cross-architecture transfer. Orbit-PE as studied in this work augments SANE's tokenization with orbit membership information derived from the channel permutation group.

ProbeGen [Kahana, Horwitz et al., 2024] achieves 30-1000x fewer FLOPs than SANE via learned probe generators. The efficiency gains from the ProbeGen line of work motivate the <= 1.2x overhead constraint adopted in H-M1.

\subsubsection{2.3 GL Orbit Symmetry and Weight Space Geometry}

Transformer-NFN [Tran-Viet et al., arXiv:2410.04209] proposes GL-invariant polynomial trace features (tr(WW^T) and tr(W^Q W^{K,T})) for attention weights, motivated by the argument that GL symmetry is the relevant group for attention weight structure. The H-M2 finding in the present work — that Linear (FC) layers show GL-orbit variance 6.6x larger than permutation orbit variance — provides independent empirical evidence from a different angle: direct variance decomposition at zoo scale on CNN Zoo linear layers, rather than theoretical argument. The Transformer Zoo attention layers were not directly measured in this work (see Limitation L3).

Loss landscape analysis and neural network trajectory studies have documented flat directions along symmetry orbits in trained networks. The present finding that the permutation variance ratio decreases from approximately 0.49 to approximately 0.28 over 50 training epochs is consistent with the broader literature on symmetry breaking during optimization.

\subsubsection{2.4 Weight Space Learning Surveys}

The WSL Survey \citep{Hanetal2026} introduces the WSU/WSR/WSG taxonomy and identifies cross-architecture generalization as a primary open problem. The present work differs from prior weight space learning studies in focusing on direct measurement of which symmetry groups explain dominant weight-space variance by layer type, rather than assuming a symmetry group and constructing an equivariant architecture.

